% A3_CIF_paper.tex
% For Physical Review X Intelligence (PRX-I) submission
% Using REVTeX 4.2 (APS format)

\documentclass[aps,prx,reprint,twocolumn,superscriptaddress,longbibliography]{revtex4-2}

\usepackage{graphicx}
\usepackage{amsmath,amssymb}
\usepackage{hyperref}
\usepackage{xcolor}
\usepackage{booktabs}

\begin{document}

\title{Backfire Phase Structure in Contaminated Chain-of-Thought Reasoning}

\author{HIDEKI}
\affiliation{Independent Researcher, Japan}

\date{\today}

\begin{abstract}
When large language models consume chain-of-thought reasoning generated by other models, errors in the consumed trace can cause the consumer to abandon its own correct reasoning---a phenomenon we term \emph{contamination-induced failure} (CIF). Here we report a systematic study of CIF across five language models, three reasoning domains, and four contamination levels (5{,}000 trials, $n=100$ per cell). We find that CIF rates vary from 1\% to 86\%, forming a phase structure that organizes along two axes: domain answer verifiability and a model-side factor consistent with what we term epistemic robustness. Crucially, this structure is not explained by sycophancy or social compliance: the most socially compliant model in our battery (Claude Sonnet~4) exhibits the lowest CIF across all domains ($\leq$7\%), while less compliant models show catastrophic vulnerability. Small differences in baseline capability produce CIF-rate ratios exceeding 57:1---a gap amplification effect with direct implications for AI pipeline safety. These findings establish CIF as a distinct failure mode requiring evaluation frameworks beyond standard accuracy benchmarks.
\end{abstract}

\maketitle

%% ============================================================
%% INTRODUCTION
%% ============================================================

\section{Introduction}

The deployment of large language models (LLMs) increasingly involves multi-agent pipelines in which the output of one model serves as input to another. Retrieval-augmented generation systems feed externally sourced reasoning into a consumer model's context. Automated tutoring platforms present AI-generated explanations to student-facing models. Code generation workflows chain multiple models in sequence, each consuming the previous model's reasoning trace. In all these settings, a consumer model must evaluate reasoning it did not produce---reasoning that may contain errors.

A natural expectation is that stronger models will be more resilient to such errors, and that this resilience will scale smoothly with model capability. Under this view, the risk of consuming flawed reasoning is a monotonically decreasing function of the consumer's competence: better models make fewer mistakes, including the mistake of trusting bad reasoning. If this were true, capability benchmarks would serve as adequate proxies for pipeline safety.

We show that this expectation is wrong in a specific and consequential way. When models consume chain-of-thought traces containing planted errors, some model-domain combinations exhibit near-total immunity (CIF rates below 4\%), while others show catastrophic vulnerability (rates above 85\%). These results do not fall along a graded spectrum. Instead, the transition between resilience and vulnerability appears abrupt, and small differences in baseline accuracy produce CIF-rate ratios exceeding 57:1---a phenomenon we term \emph{gap amplification}.

The most surprising feature of this phase structure is what does \emph{not} explain it. The sycophancy literature~\cite{sharma2023towards,perez2023discovering,wei2023simple} has characterized how models trained with human feedback tend to defer to user opinions rather than maintain truthful responses. A natural extension of this framework to multi-model pipelines would predict that contamination-induced failure is driven by the same mechanism: models fail because they comply with the stated answer in the consumed trace. We show this prediction is incorrect. Claude Sonnet~4 exhibits the highest measured social compliance---a cue-override rate of 88\% on our E3 battery---yet the lowest CIF across all domains ($\leq$7\%). GPT-4o-mini, with a cue-override rate of only 15\%, shows CIF rates up to 34\%. Social compliance and contamination vulnerability are dissociated.

This dissociation points toward an alternative mechanism. We propose that CIF is modulated by \emph{epistemic robustness}: the capacity to internally verify or reconstruct the reasoning steps in a consumed trace, independent of social pressure to agree. On this account, the relevant axis is not whether a model tends to agree, but whether it can \emph{check}. This reframing explains the observed domain gradient---mathematical reasoning (GSM8K), where answers are step-by-step verifiable, shows lower CIF than commonsense reasoning (CSQA) or reading comprehension (BoolQ), where verification requires holistic judgment---and the model ordering, in which capability differences are amplified into qualitatively different failure regimes.

The contribution of this paper is empirical and conceptual. Empirically, we establish the phenomenon: contamination-induced failure is real, large, and structured. We provide a controlled experimental framework (the CIF protocol) that isolates the effect of trace contamination from confounds such as problem difficulty or prompt phrasing. Conceptually, we argue that standard accuracy benchmarks and sycophancy evaluations are insufficient for predicting pipeline safety, and that a two-dimensional characterization---model robustness $\times$ domain verifiability---is necessary to identify hazardous operating regimes.

We do not claim to have identified the mechanistic basis of epistemic robustness, nor do we propose a predictive law. What we report is a phenomenon that existing frameworks fail to explain and that has immediate practical consequences for the design and evaluation of multi-model systems.


%% ============================================================
%% EXPERIMENTAL FRAMEWORK
%% ============================================================

\section{Experimental Framework}

\subsection{The CIF Protocol}

We designed a controlled protocol to measure how consuming contaminated chain-of-thought reasoning affects a model's accuracy on problems it can solve independently. The protocol involves two roles: a \emph{producer} model that generates reasoning traces, and a \emph{consumer} model that solves problems while presented with those traces.

\textbf{Trace generation.}
For each problem, GPT-4o (the producer) generates two 8-step reasoning traces: a \emph{clean trace} leading to the correct answer and a \emph{corrupt trace} leading to a specific wrong answer. The corrupt trace contains plausible but erroneous reasoning---not random noise, but the kind of mistake a competent reasoner might make.

\textbf{Contamination mixing.}
A mixing parameter $\lambda \in \{0.0, 0.4, 0.8, 1.0\}$ controls the fraction of corrupt steps substituted into the trace. At $\lambda = 0$, the consumer sees a fully correct trace. At $\lambda = 1$, every step is corrupt. Intermediate values produce traces where some steps are correct and others are not---a realistic model of partially reliable external reasoning. Step selection is deterministic given a seed, ensuring reproducibility.

\textbf{Consumer evaluation.}
Each consumer model solves each problem under two conditions: (1)~\emph{direct}, with no reasoning trace provided, establishing a baseline accuracy; and (2)~\emph{wrong-cue}, with the mixed trace and a stated (incorrect) answer appended to the prompt. The direct baseline prompt requires JSON-only output without an explicit reasoning trace. While frontier models appear able to maintain internal reasoning under this constraint, we observe that some mid-tier models show reduced baseline accuracy on harder problems---an effect we attribute to prompt-induced suppression of chain-of-thought reasoning rather than model instability (see Supplementary Fig.~S1). Temperature is fixed at~0 for all consumer calls.

\textbf{CIF metric.}
Contamination-induced failure is defined as
\begin{equation}
\mathrm{CIF}(m, d, \lambda) = P(\text{incorrect} \mid \text{wrong-cue}, \lambda, \text{correct at baseline}),
\label{eq:cif}
\end{equation}
where $m$ indexes the consumer model and $d$ the reasoning domain. CIF measures the fraction of problems that a model gets right on its own but gets wrong when presented with a contaminated trace. This conditional definition isolates the effect of contamination from baseline capability differences. Because CIF is computed only over problems solved correctly at baseline, reductions in baseline accuracy reduce the eligible sample size for CIF estimation rather than inflating the CIF rate itself.

\subsection{Models}

Five consumer models spanning three capability tiers: Claude Sonnet~4 and GPT-4o (frontier), Claude Haiku~3.5 and GPT-4o-mini (mid-tier), and GPT-3.5-turbo (legacy). All accessed via API with temperature~0 and maximum output tokens of 1024.

\subsection{Domains}

Three reasoning domains chosen to span a \emph{verifiability gradient}:

\emph{GSM8K}~\cite{cobbe2021training} (grade-school mathematics): answers are numerical, each step is independently checkable, and errors propagate deterministically. High verifiability.

\emph{CSQA}~\cite{talmor2019commonsenseqa} (CommonsenseQA): answers require commonsense inference across multiple plausible options. Medium verifiability.

\emph{BoolQ}~\cite{clark2019boolq} (Boolean Questions): answers are True/False based on passage comprehension, with limited step-by-step decomposability. Low verifiability.

This ordering reflects the degree to which a consumer model can internally reconstruct and verify the reasoning in a consumed trace, independent of whether it ``agrees'' with the stated conclusion.

\subsection{Social Compliance Battery (E3)}

To test whether sycophancy explains CIF, we administered a separate compliance battery (E3) measuring three dimensions: (1)~\emph{cue override rate}---how often a model switches from its own correct answer when presented with an authoritative but wrong alternative; (2)~\emph{authority conflict rate}---susceptibility to appeals to authority; and (3)~\emph{instruction resistance}---refusal to follow explicitly wrong instructions. These are combined into a composite compliance index ranging from~0 (fully independent) to~1 (fully compliant). Note that cue override rate and composite index measure different aspects: a model may show high cue override (readily updating answers when challenged) while scoring moderately on the composite (which also weights authority and instruction resistance).

\subsection{Scale}

The main experiment comprises two rounds (E1/E1b at $n=50$, E1aug at $n=50$) yielding $n=100$ problems per domain across $5$~models $\times$ $3$~domains $\times$ $5$~conditions ($1$~direct $+$ $4$~$\lambda$ levels), totaling approximately 8{,}750 consumer evaluations. The compliance battery (E3) adds 1{,}000 trials across the five models. All results reported below use the merged $n=100$ dataset unless otherwise noted.


%% ============================================================
%% RESULTS
%% ============================================================

\section{Results}

\subsection{The CIF Phase Diagram}

Figure~\ref{fig:main} reveals a phase structure in contamination-induced failure, separating models into immune, transition, and vulnerable regimes across reasoning domains. Peak CIF rates across the 15 model-domain cells range from 1.1\% (Claude Sonnet~4 on BoolQ) to 85.5\% (GPT-3.5 on CSQA), spanning more than a 75-fold range.

\textbf{Immune regime.}
Claude Sonnet~4 shows peak CIF $\leq$6.6\% across all three domains (95\% CI upper bounds below 14\%). Contaminated traces---even fully corrupt ones---fail to reliably override its reasoning. The domain gradient is effectively flat: verifiability makes little difference when the model can resist contamination regardless.

\textbf{Vulnerable regime.}
GPT-3.5 and GPT-4o-mini occupy the opposite corner. GPT-3.5 reaches 85.5\% CIF on CSQA [76\%, 92\%] and 63.5\% on BoolQ [53\%, 73\%]; GPT-4o-mini reaches 33.8\% on CSQA [24\%, 45\%] and 32.6\% on BoolQ [24\%, 43\%]. For these models, consuming a contaminated trace is worse than not reasoning at all.

\textbf{Transition zone.}
GPT-4o and Claude Haiku~3.5 fall between: partially protected in high-verifiability domains (GSM8K CIF of 3.2\% and 16.5\%, respectively) but increasingly vulnerable as verifiability decreases (BoolQ CIF of 14.9\% and 15.2\%).

Quantitatively, the model ordering by mean peak CIF is: Sonnet~4 (3.6\%) $<$ GPT-4o (11.6\%) $<$ Haiku~3.5 (15.4\%) $<$ GPT-4o-mini (26.9\%) $<$ GPT-3.5 (54.7\%). The domain ordering follows the verifiability gradient: GSM8K (9.9\%) $<$ CSQA (26.9\%) $<$ BoolQ (24.7\%). Critically, the two axes interact: domain effects are negligible for immune models (Sonnet~4 varies by less than~6 percentage points across domains) but dominate for vulnerable ones (GPT-3.5 ranges from 15\% to 86\%). The phase structure is not a simple additive decomposition of model and domain effects.

\subsection{Lambda Response Curves}

The contamination parameter $\lambda$ modulates CIF nonlinearly. For vulnerable models, the transition is steep: GPT-3.5 on CSQA jumps from 5\% CIF at $\lambda=0$ to 63\% at $\lambda=0.4$, with diminishing increases thereafter (79\% at $\lambda=0.8$, 82\% at $\lambda=1.0$). This pattern---a sharp onset followed by saturation---is characteristic of a threshold phenomenon rather than proportional degradation.

For robust models, the curve is flat. Claude Sonnet~4 on GSM8K shows approximately 3--4\% CIF at every $\lambda$ level, indicating that its resistance does not depend on the fraction of corrupt steps---even a fully corrupt trace fails to reliably mislead it.

An intermediate pattern appears for mid-capability models in high-verifiability domains: GPT-4o-mini on GSM8K rises from 2\% ($\lambda=0$) to 21\% ($\lambda=0.4$), then partially retracts to 15\% ($\lambda=0.8$) and 17\% ($\lambda=1.0$), suggesting that fully corrupt traces may be easier to detect than partially corrupt ones.

\subsection{Gap Amplification}

The most practically significant finding is the nonlinear relationship between baseline capability and CIF. We define the \emph{gap amplification factor} (GAF) as the ratio of peak CIF rates between two models on the same domain:
\begin{equation}
\mathrm{GAF}(A, B, d) = \frac{\max_\lambda \mathrm{CIF}(A, d, \lambda)}{\max_\lambda \mathrm{CIF}(B, d, \lambda)}.
\label{eq:gaf}
\end{equation}
On GSM8K, where baseline accuracies are tightly clustered among frontier and mid-tier models, GAF between Sonnet~4 and Haiku~3.5 is 5.4$\times$. On CSQA, the Sonnet~4\,/\,GPT-3.5 GAF reaches 13$\times$. On BoolQ, it reaches 57$\times$---a model with only moderately lower baseline accuracy is an order of magnitude more vulnerable to contaminated reasoning, a relationship that no accuracy benchmark would predict (Fig.~\ref{fig:main}C).

\subsection{Sycophancy Does Not Explain CIF}

The compliance battery (E3) reveals a striking dissociation. Claude Sonnet~4 exhibits the highest cue-override rate among all models tested (88\%---it switches its answer when presented with an authoritative alternative on nearly nine out of ten eligible problems) and the highest composite compliance index among frontier models (0.29). By standard sycophancy measures, it is the most deferential model in the set. Yet it shows the lowest CIF ($\leq$6.6\%). Conversely, GPT-4o-mini has the lowest cue-override rate (15\%) and the lowest composite compliance index (0.05) but CIF rates up to 34\%.

If sycophancy (social compliance with stated answers) were the mechanism driving CIF, compliance index and CIF should correlate positively. They do not. The rank ordering is inconsistent, and the most compliant model is the most resistant to contamination.

This falsification is the conceptual core of the paper. The sycophancy framework predicts that models fail under contaminated input because they defer to the stated answer. Our data show that deferral and failure are dissociated: a model can be socially compliant (willing to update its answer when challenged) while remaining epistemically robust (able to recognize when a reasoning trace is wrong). These are different capacities.


%% ============================================================
%% DISCUSSION
%% ============================================================

\section{Discussion}

\subsection{Epistemic Robustness as a Candidate Mechanism}

The phase structure itself---three regimes, two axes, gap amplification---is an empirical finding. What \emph{explains} this structure is a separate question. The dissociation between social compliance and contamination vulnerability suggests a candidate: \emph{epistemic robustness}, the capacity to internally verify the logical coherence of consumed reasoning, independent of social cues. A model with high epistemic robustness can follow a chain-of-thought, detect inconsistencies or errors within it, and override the stated conclusion when its own verification fails. We propose this as a hypothesis, not a demonstrated mechanism.

This hypothesis explains three features of our data. First, the domain gradient: mathematical reasoning is more internally verifiable (each arithmetic step can be checked) than commonsense or passage-based reasoning, so epistemic robustness would protect more effectively in GSM8K than in BoolQ. Second, the model ordering: frontier models with stronger internal representations may perform deeper verification, producing lower CIF even when they are socially compliant. Third, gap amplification: if verification capacity exhibits threshold behavior---below a certain capability level, models cannot effectively check consumed reasoning, and above it, they can---then small capability differences near this threshold would translate to large differences in verification success.

We emphasize the distinction: the \emph{phenomenon} (CIF phase structure, gap amplification, sycophancy dissociation) is established by our data. The \emph{mechanism} (epistemic robustness as the governing variable) is a hypothesis consistent with the data but requiring independent validation---through probing internal representations, attention pattern analysis, or targeted ablation.

\subsection{Relation to Prior Work}

Our findings connect to several research streams while identifying gaps in existing frameworks.

\emph{Chain-of-thought reasoning.}
Chain-of-thought prompting~\cite{wei2022chain} and self-consistency decoding~\cite{wang2023self} have established that explicit reasoning traces improve LLM performance. However, this body of work focuses on model-generated reasoning. Our protocol inverts the question: what happens when a model consumes reasoning it did not generate? The benefits of chain-of-thought become a liability when the consumed trace contains errors.

\emph{CoT faithfulness.}
Work on CoT faithfulness~\cite{turpin2023language,lanham2023measuring} has shown that models' stated reasoning does not always reflect their actual decision process---biased inputs can shift CoT explanations without the model acknowledging the influence. More broadly, the question of whether model explanations faithfully reflect internal processes has been a concern since before the era of large language models~\cite{jacovi2020towards}. Turpin~et~al.\ study faithfulness of internally generated reasoning, while we study robustness to externally provided reasoning. The two questions are orthogonal but may share an underlying mechanism---a model that generates faithful CoT may also be better at evaluating others'---though this connection remains to be tested. Recent work on process-based verification~\cite{lightman2023lets} demonstrates that step-by-step verification can improve reasoning quality, providing indirect support for our hypothesis that verification capacity underlies the domain gradient in CIF.

\emph{Sycophancy research.}
The sycophancy literature~\cite{sharma2023towards,perez2023discovering,wei2023simple} has established that LLMs often defer to user opinions, but has primarily studied this in conversational and opinion-based contexts. Our results suggest that the sycophancy construct may conflate two distinct phenomena: social compliance (updating answers under social pressure) and epistemic capitulation (abandoning correct reasoning under evidential pressure). Concurrent work has shown that sycophancy itself is not a unitary construct---sycophantic agreement and sycophantic praise are causally separable behaviors~\cite{kim2025sycophancy}. Our results extend this insight: contamination vulnerability (CIF) is dissociated not only from sycophantic agreement but from social compliance more broadly.

\emph{Multi-agent system failures.}
Recent systematic analysis of multi-agent LLM failures reports that task verification failures account for roughly one-fifth of observed system breakdowns~\cite{cemri2025multi}. Multi-agent debate frameworks~\cite{du2023improving} assume that inter-model reasoning exchange improves accuracy, but our findings reveal conditions under which the opposite occurs: consuming another model's flawed reasoning can be worse than reasoning independently. The CIF protocol provides a controlled method for measuring this vulnerability before deployment, complementing the post-hoc failure taxonomy of Cemri~et~al.

\emph{LLM-as-judge.}
The growing practice of using LLMs to evaluate other LLMs' outputs~\cite{zheng2023judging} assumes that the evaluating model can reliably detect errors. Our CIF measurements provide a direct empirical test of this assumption, and the results are cautionary: mid-tier and legacy models show substantial failure rates when consuming flawed reasoning, suggesting that judge model selection is a safety-critical decision. None of these prior works examine how models evaluate externally provided reasoning traces under controlled contamination conditions---the specific gap that the CIF protocol addresses.

\subsection{Practical Implications}

The gap amplification effect has immediate consequences for AI system design. In a pipeline where a weaker model consumes reasoning from a stronger model that has made an error, the weaker model's vulnerability is not proportional to its capability deficit---it is amplified. On BoolQ, a model whose baseline accuracy differs only moderately from Sonnet~4's shows a CIF rate 57~times higher. This means that standard accuracy benchmarks systematically underestimate the risk of multi-model pipelines.

We recommend three practices based on our findings: (1)~evaluate pipeline components under contaminated-input conditions, not only clean benchmarks---extending the behavioral testing philosophy of Ribeiro~et~al.~\cite{ribeiro2020beyond} to multi-model settings; (2)~match consumer model capability to the verifiability of the reasoning domain---high-verifiability tasks tolerate weaker consumers, but low-verifiability tasks require frontier-tier consumers; (3)~treat the CIF phase diagram as a risk map, identifying model-domain cells where contamination vulnerability is unacceptably high.


%% ============================================================
%% LIMITATIONS
%% ============================================================

\section{Limitations}

Several limitations qualify the scope of our findings. First, our sample size of 100~problems per domain yields 95\% Wilson confidence intervals of approximately $\pm$8--12 percentage points for CIF rates in the 15--40\% range. The qualitative phase ordering---immune, transition, vulnerable---is stable across independent $n=50$ sweeps, but individual cell estimates should be interpreted with their reported confidence intervals (Fig.~\ref{fig:main}B). Second, three reasoning domains, while chosen to span a verifiability gradient, do not constitute a dense sampling of the verifiability axis. Additional domains (e.g., formal logic, causal reasoning, spatial reasoning) would strengthen the gradient claim. Third, the epistemic robustness hypothesis is supported by the pattern of results but lacks direct mechanistic evidence. We have measured the \emph{consequences} of differential robustness without probing the \emph{process} by which models verify or fail to verify consumed reasoning. Fourth, all experiments use a single producer model (GPT-4o) for trace generation, and producer-consumer interactions may differ with other producers. Fifth, our protocol uses explicit chain-of-thought presentation; implicit contamination through retrieval-augmented generation or fine-tuning may produce different vulnerability profiles. Sixth, the direct baseline prompt suppresses externalized chain-of-thought (JSON-only output), which reduces baseline accuracy for some mid-tier models on harder problems (see Supplementary Fig.~S1). Because CIF is conditional on baseline correctness, this effect reduces the eligible sample size for those cells rather than inflating CIF rates, but it widens confidence intervals for the affected cells. Finally, our E3 compliance battery measures social compliance under explicit challenge conditions, which may not capture all dimensions of the sycophancy construct identified by recent work~\cite{sharma2023towards,kim2025sycophancy}. In particular, argument-driven sycophancy---updating beliefs under reasoned argumentation---may operate through different mechanisms than cue-override sycophancy.


%% ============================================================
%% CONCLUSION
%% ============================================================

\section{Conclusion}

We have reported the discovery of a structured pattern in how language models fail when consuming contaminated chain-of-thought reasoning. Across five models, three domains, and four contamination levels ($n=100$ per cell), contamination-induced failure rates range from 1\% to 86\%, organizing into a phase structure along two axes: model epistemic robustness and domain answer verifiability.

Three findings carry particular weight. First, gap amplification: small differences in baseline capability produce CIF-rate ratios exceeding 57:1, meaning that standard benchmarks dramatically underestimate the differential risk in multi-model pipelines. Second, the sycophancy falsification: the model with the highest social compliance shows the lowest contamination vulnerability, demonstrating that CIF is not a special case of sycophancy but a distinct failure mode. Third, the domain gradient: mathematical reasoning with step-by-step verifiability is systematically more resistant than passage-based comprehension, pointing to internal verification capacity as a candidate mechanism.

These results imply that evaluating LLMs solely on accuracy or sycophancy benchmarks leaves a critical safety dimension unmeasured. As AI systems increasingly consume each other's reasoning---through retrieval, cascading, and multi-agent architectures---the CIF phase structure becomes a risk map. Without this framework, deploying a model safe on math benchmarks into a RAG pipeline for commonsense reasoning risks catastrophic failure rates an order of magnitude higher than benchmarks predict---and sycophancy evaluations will not catch this. Identifying and characterizing these failure regimes before deployment is not optional; it is a prerequisite for trustworthy AI systems.


%% ============================================================
%% DATA AVAILABILITY
%% ============================================================

\section*{Data Availability}

All experimental data, analysis code, and reproducibility materials are available at \url{https://github.com/HIDEKI-SQ/Contamination-Induced-Failure} [to be updated with DOI upon archival].


%% ============================================================
%% ACKNOWLEDGMENTS
%% ============================================================

\begin{acknowledgments}
The author acknowledges the use of AI-assisted tools for feedback on manuscript structure and for supporting data analysis workflows. All scientific claims, experimental decisions, and interpretations are solely the author's responsibility.
\end{acknowledgments}


%% ============================================================
%% FIGURES
%% ============================================================

\begin{figure*}[t]
\centering
\includegraphics[width=\textwidth]{Figure1_n100.pdf}
\caption{\textbf{Contamination-induced failure phase structure ($n=100$).}
\textbf{(A)}~CIF rate as a function of contamination level~$\lambda$ for five models across three reasoning domains (GSM8K, CSQA, BoolQ). Error bars show 95\% Wilson confidence intervals. Colored bands indicate baseline accuracy for each model (right axis, light shading). Sonnet~4 (blue) remains flat near zero across all conditions; GPT-3.5 (red) saturates at high CIF after a sharp onset at $\lambda=0.4$.
\textbf{(B)}~CIF phase diagram showing peak CIF with 95\% CI for each model-domain cell. Background shading delineates three regimes: immune (green, $<$10\%), transition (orange, 10--40\%), and vulnerable (red, $>$40\%). Point opacity and $n$-labels indicate eligible sample size---cells with lower baseline accuracy have fewer eligible problems. Sonnet~4 clusters in the immune regime across all domains; GPT-3.5 occupies the vulnerable regime, particularly for low-verifiability tasks.
\textbf{(C)}~Gap amplification factor (GAF\,$=$\,peak CIF ratio vs.\ Sonnet~4) by domain. GAF increases from $\sim$5$\times$ on GSM8K to 57$\times$ on BoolQ for GPT-3.5, demonstrating that capability differences are amplified as domain verifiability decreases.}
\label{fig:main}
\end{figure*}


%% ============================================================
%% REFERENCES
%% ============================================================

\bibliography{references}

\begin{thebibliography}{13}

\bibitem{sharma2023towards}
M.~Sharma, M.~Tong, T.~Korbak, D.~Duvenaud, A.~Askell, S.~R.~Bowman, \emph{et~al.},
Towards understanding sycophancy in language models,
arXiv:2310.13548 (2023).

\bibitem{perez2023discovering}
E.~Perez, S.~Ringer, K.~Luko\v{s}i\=ut\.e, K.~Nguyen, \emph{et~al.},
Discovering language model behaviors with model-written evaluations,
in \emph{Findings of ACL 2023}, pp.~13387--13434 (2023).

\bibitem{wei2023simple}
J.~Wei, D.~Huang, Y.~Lu, D.~Zhou, and Q.~V.~Le,
Simple synthetic data reduces sycophancy in large language models,
arXiv:2308.03958 (2023).

\bibitem{cobbe2021training}
K.~Cobbe, V.~Kosaraju, M.~Bavarian, M.~Chen, H.~Jun, L.~Kaiser, \emph{et~al.},
Training verifiers to solve math word problems,
arXiv:2110.14168 (2021).

\bibitem{talmor2019commonsenseqa}
A.~Talmor, J.~Herzig, N.~Lourie, and J.~Berant,
CommonsenseQA: A question answering challenge targeting commonsense knowledge,
in \emph{Proc.\ NAACL-HLT 2019}, pp.~4149--4158 (2019).

\bibitem{clark2019boolq}
C.~Clark, K.~Lee, M.-W.~Chang, T.~Kwiatkowski, M.~Collins, and K.~Toutanova,
BoolQ: Exploring the surprising difficulty of natural yes/no questions,
in \emph{Proc.\ NAACL-HLT 2019}, pp.~2924--2936 (2019).

\bibitem{wei2022chain}
J.~Wei, X.~Wang, D.~Schuurmans, M.~Bosma, B.~Ichter, F.~Xia, E.~Chi, Q.~Le, and D.~Zhou,
Chain-of-thought prompting elicits reasoning in large language models,
in \emph{NeurIPS~35} (2022).

\bibitem{wang2023self}
X.~Wang, J.~Wei, D.~Schuurmans, Q.~Le, E.~Chi, S.~Narang, A.~Chowdhery, and D.~Zhou,
Self-consistency improves chain of thought reasoning in language models,
arXiv:2203.11171 (2023).

\bibitem{turpin2023language}
M.~Turpin, J.~Michael, E.~Perez, and S.~R.~Bowman,
Language models don't always say what they think: Unfaithful explanations in chain-of-thought prompting,
in \emph{NeurIPS~36} (2023).

\bibitem{lanham2023measuring}
T.~Lanham, A.~Chen, A.~Radhakrishnan, J.~Stehardt, S.~Bowman, \emph{et~al.},
Measuring faithfulness in chain-of-thought reasoning,
arXiv:2307.13702 (2023).

\bibitem{kim2025sycophancy}
S.~Kim \emph{et~al.},
Sycophancy is not one thing: Causal separation of sycophantic agreement and sycophantic praise,
arXiv:2509.21305 (2025).

\bibitem{cemri2025multi}
M.~Cemri, M.~Z.~Pan, S.~Yang, \emph{et~al.},
Why do multi-agent LLM systems fail?,
arXiv:2503.13657 (2025).

\bibitem{zheng2023judging}
L.~Zheng, W.-L.~Chiang, Y.~Sheng, S.~Zhuang, Z.~Wu, \emph{et~al.},
Judging LLM-as-a-judge with MT-Bench and Chatbot Arena,
in \emph{NeurIPS~36} (2023).

\bibitem{jacovi2020towards}
A.~Jacovi and Y.~Goldberg,
Towards faithfully interpretable NLP systems: A survey,
in \emph{Proc.\ ACL 2020}, pp.~4198--4205 (2020).

\bibitem{ribeiro2020beyond}
M.~T.~Ribeiro, T.~Wu, C.~Guestrin, and S.~Singh,
Beyond accuracy: Behavioral testing of NLP models with CheckList,
in \emph{Proc.\ ACL 2020}, pp.~4902--4912 (2020).

\bibitem{lightman2023lets}
H.~Lightman, V.~Kosaraju, Y.~Burda, H.~Edwards, B.~Baker, T.~Lee, J.~Leike, J.~Schulman, I.~Sutskever, and K.~Cobbe,
Let's verify step by step,
arXiv:2305.20050 (2023).

\bibitem{du2023improving}
Y.~Du, S.~Li, A.~Torralba, J.~B.~Tenenbaum, and I.~Mordatch,
Improving factuality and reasoning in language models through multiagent debate,
arXiv:2305.14325 (2023).

\end{thebibliography}

\end{document}
